\chapter{Resultados}
\label{ch:resultados}
% Única fuente numérica: docs/evaluation/FINAL_HOLDOUT_V2_RESULTS.md.
% Copiar las tablas A--F sin ejecutar ni modificar el evaluador.
\section{Composición y resultado operativo}
\pendiente{Transferir la tabla A y los estados de la sección 7 del documento
canónico: 47 TERMINAL\_SUCCESS, un TERMINAL\_ERROR y ningún pendiente o
indeterminado. Éxito operativo no significa corrección semántica.}
\section{Resultados primarios: clasificación documental}
\pendiente{PRIMARY RESULTS. Transferir alcance documental y matriz D,
con su denominador completo y la predicción ausente.}
\section{Resultados primarios: detección de entidades}
\pendiente{PRIMARY RESULTS. Transferir tabla B: activos, eventos,
actuaciones actuales y localizaciones, en ese orden. Mantener TP, FP, FN,
precisión, recall y F1 y los denominadores originales.}
\section{Resultados primarios: atributos condicionados}
\pendiente{PRIMARY RESULTS. Transferir tabla C y explicar a qué pares
aplicables se limita cada atributo. Un acierto sobre un caso no demuestra
rendimiento general.}
\section{Resultados primarios: actuación y activo afectado}
\pendiente{PRIMARY RESULTS. Transferir tabla E, unidad relacional y
condicionamiento por emparejamiento previo; no es evaluación de grouping.}
\section{Resultados primarios: evidencia y P0}
\pendiente{PRIMARY RESULTS. Transferir tabla F; conservar denominadores
pequeños y valores no definidos. No convertirlos en ceros ni completar con
resultados del desarrollo. El diagnóstico posterior se presenta por separado
en el capítulo~\ref{ch:discusion}.}
